% fixture: babel/spanish-basic
% tier: T3
% engine: document
% covers: \usepackage[spanish]{babel} (spanish.ldf 5.0q, defaults): \captionsspanish (Resumen, \'{I}ndice, Figura, Cuadro, Referencias), \spanishdate's \today in \maketitle, \frenchspacing, the spanish patterns with hyphenmins 2 2, \es@accent's \bbl@allowhyphens after OT1 \' and \~ unless a space follows, \layoutspanish (numbers with a period in headings and the contents, the first paragraph after a heading indented), \% with a thin space
% issue: #63
\documentclass{article}
\usepackage[spanish]{babel}
\title{Título del artículo}
\author{José Pérez}
\begin{document}
\maketitle
\begin{abstract}
Un resumen corto del artículo. ¿Qué tal? ¡Bien! El Dr. Pérez dijo: sí. Fin.
\end{abstract}
\tableofcontents
\section{Introducción}\label{s}
La teoría de la información estudia la cuantificación, el almacenamiento y la comunicación de la
información digital. Fue establecida por Claude Shannon en los años cuarenta. Véase la
sección~\ref{s}, la figura~\ref{f} y el cuadro~\ref{t}.\footnote{Una nota con acentuación.} Más texto:
niño, año, \'arbol, ma\~nana, \'{\i}ndice y \'{I}ndice; el 50\% y el 50 \% del total, \'e.
\begin{figure}[h]
\caption{Descripción del método}\label{f}
\end{figure}
\begin{table}[h]
\caption{Resultados numéricos}\label{t}
\end{table}
\begin{minipage}{60pt}
Transformaciones establecidas y almacenamiento de comunicaciones electrónicas.
\end{minipage}
\subsection{Conclusión}
Fin del artículo.
\begin{thebibliography}{9}
\bibitem{a} Un ítem.
\end{thebibliography}
\end{document}
